
\chapter{Poincaré Duality}
\label{ch:viii33-poincare-duality}

Poincaré duality is the organizing theorem that turns cohomology classes into
geometric homology classes on an oriented manifold.  The Thom-class viewpoint
from VIII/31 supplies the local normal model; the Euler-class viewpoint from
VIII/32 supplies a first family of dual cycles.

\section{Orientation and local homology}
For an $n$-manifold $M$ and $x\in M$,
\[
H_n(M,M\setminus\{x\};\mathbb Z)\cong\mathbb Z.
\]
An orientation chooses compatible generators of these local groups.

\section{The fundamental class}
\begin{definition}[Fundamental class]
\label{def:viii33-fundamental-class}
For a connected closed oriented $n$-manifold, the fundamental class
\[
[M]\in H_n(M;\mathbb Z)
\]
is the unique class restricting to the chosen local orientation generator at
every point.
\end{definition}

\section{Cap product}
Cap product combines cohomology and homology:
\[
\frown:H^k(M;R)\times H_n(M;R)\longrightarrow H_{n-k}(M;R).
\]

\section{Poincaré duality}
\begin{theorem}[Poincaré duality]
\label{thm:viii33-poincare-duality}
If $M$ is a closed oriented $n$-manifold, then
\[
D_M:H^k(M;R)\longrightarrow H_{n-k}(M;R),
\qquad
D_M(\alpha)=\alpha\frown[M]
\]
is an isomorphism.
\end{theorem}

\section{Betti-number symmetry}
Over a field,
\[
b_k(M)=b_{n-k}(M).
\]
This gives immediate global restrictions, including vanishing Euler
characteristic in odd dimension.

\section{The complementary-degree pairing}
Duality can be written as the nondegenerate pairing
\[
H^k(M;R)\times H^{n-k}(M;R)\longrightarrow R,
\qquad
(\alpha,\beta)\longmapsto
\langle\alpha\smile\beta,[M]\rangle.
\]

\section{Spheres and surfaces}
For $S^n$, degree zero and degree $n$ are dual.  For a closed oriented surface
of genus $g$, the $2g$ degree-one generators pair through the symplectic
intersection form.

\section{The torus}
Choose $a,b\in H^1(T^2;\mathbb Z)$ so that
\[
a\smile b=[T^2]^*,\qquad b\smile a=-[T^2]^*.
\]
The algebra records the signed intersection of the two standard generating
circles.

\section{Complex projective space}
For $\mathbb{CP}^n$,
\[
H^*(\mathbb{CP}^n;\mathbb Z)\cong\mathbb Z[h]/(h^{n+1}),
\qquad |h|=2,
\]
and the top class $h^n$ evaluates to $1$ after choosing the complex
orientation.

\section{Geometric Poincaré duals}
If $N^{n-r}\hookrightarrow M^n$ is a closed oriented submanifold with oriented
normal bundle, its normal Thom class extends to
\[
\operatorname{PD}[N]\in H^r(M).
\]

\section{Transverse intersection and cup product}
If dual submanifolds meet transversely, then cup product corresponds to
intersection:
\[
\operatorname{PD}[A]\smile\operatorname{PD}[B]
=
\operatorname{PD}[A\pitchfork B]
\]
with the usual orientation conventions.

\section{Intersection numbers}
For complementary dimensions, the signed number of transverse intersection
points is
\[
A\cdot B
=
\langle \operatorname{PD}[A]\smile\operatorname{PD}[B],[M]\rangle.
\]

\section{Euler classes revisited}
If a transverse section of an oriented rank-$r$ bundle has zero set $Z$, then
\[
\operatorname{PD}[Z]=e(E).
\]
Thus Euler classes are geometric duals of zero loci.

\section{Manifolds with boundary}
\begin{theorem}[Poincaré--Lefschetz duality]
\label{thm:viii33-poincare-lefschetz}
For a compact oriented $n$-manifold with boundary,
\[
H^k(M;R)\cong H_{n-k}(M,\partial M;R),
\qquad
H^k(M,\partial M;R)\cong H_{n-k}(M;R).
\]
\end{theorem}

\section{Nonorientable manifolds}
Integral duality requires the orientation local system.  With
$\mathbb Z/2$ coefficients the sign obstruction disappears and a mod-$2$
duality statement holds without orientability.

\section{Bridge to intersection forms}
In dimension $4m$, the middle-degree pairing
\[
H^{2m}(M)\times H^{2m}(M)\to R
\]
is a bilinear form obtained by cup product and evaluation on $[M]$.  VIII/34
will study this intersection form as an invariant in its own right.

\section{Solved dossiers}
\begin{problem}
State Poincaré duality for a closed oriented $n$-manifold.
\end{problem}
\begin{solution}
Cap product with $[M]$ gives $H^k(M;R)\cong H_{n-k}(M;R)$.
\end{solution}

\begin{problem}
Define the duality map.
\end{problem}
\begin{solution}
$D_M(\alpha)=\alpha\frown[M]$.
\end{solution}

\begin{problem}
What is the fundamental class?
\end{problem}
\begin{solution}
The homology class whose local image is the chosen orientation generator at every point.
\end{solution}

\begin{problem}
What symmetry follows for Betti numbers over a field?
\end{problem}
\begin{solution}
$b_k=b_{n-k}$.
\end{solution}

\begin{problem}
Write the degree-pairing form.
\end{problem}
\begin{solution}
$(\alpha,\beta)\mapsto\langle\alpha\smile\beta,[M]\rangle$.
\end{solution}

\begin{problem}
Why is the pairing nondegenerate?
\end{problem}
\begin{solution}
Poincaré duality identifies one factor with the dual of complementary homology/cohomology.
\end{solution}

\begin{problem}
Compute the dual of the point class on a connected oriented surface.
\end{problem}
\begin{solution}
It corresponds to the top-degree cohomology generator.
\end{solution}

\begin{problem}
Describe duality on $S^n$.
\end{problem}
\begin{solution}
$H^0$ pairs with $H^n$ and vice versa; all intermediate groups vanish.
\end{solution}

\begin{problem}
Describe the first-cohomology pairing on $T^2$.
\end{problem}
\begin{solution}
Two degree-one generators have nonzero cup product equal to the orientation class.
\end{solution}

\begin{problem}
How does a submanifold define a dual class?
\end{problem}
\begin{solution}
Use the Thom class of its oriented normal bundle and extend it to ambient cohomology.
\end{solution}

\begin{problem}
What does cup product represent geometrically for transverse dual submanifolds?
\end{problem}
\begin{solution}
Their intersection.
\end{solution}

\begin{problem}
What is the intersection number of complementary oriented transverse submanifolds?
\end{problem}
\begin{solution}
The evaluation of the cup product of their dual classes on $[M]$.
\end{solution}

\begin{problem}
How does boundary change the theorem?
\end{problem}
\begin{solution}
Poincaré--Lefschetz duality relates absolute cohomology to relative homology and relative cohomology to absolute homology.
\end{solution}

\begin{problem}
State one Poincaré--Lefschetz isomorphism.
\end{problem}
\begin{solution}
$H^k(M)\cong H_{n-k}(M,\partial M)$ for compact oriented $M$.
\end{solution}

\begin{problem}
What changes for nonorientable manifolds?
\end{problem}
\begin{solution}
Use the orientation local system, or mod-$2$ coefficients when appropriate.
\end{solution}

\begin{problem}
How is the Euler characteristic constrained in odd dimension?
\end{problem}
\begin{solution}
Betti numbers pair in complementary degrees, giving zero Euler characteristic for closed oriented odd-dimensional manifolds.
\end{solution}

\begin{problem}
What is the dual of a codimension-$r$ transverse zero locus?
\end{problem}
\begin{solution}
The Euler class of the rank-$r$ bundle whose section cuts it out.
\end{solution}

\begin{problem}
How is the Thom class used in the proof philosophy?
\end{problem}
\begin{solution}
Local duality in normal disks is transported globally through tubular neighborhoods and excision.
\end{solution}

\begin{problem}
Why does VIII/33 lead directly to intersection forms?
\end{problem}
\begin{solution}
Middle-degree cup products evaluated on $[M]$ become bilinear forms.
\end{solution}

\begin{problem}
What is the key synthesis of VIII/31--VIII/33?
\end{problem}
\begin{solution}
Thom classes create dual cohomology classes, Euler classes arise from zero sections, and cap product with the fundamental class identifies cohomology with geometric homology.
\end{solution}

\section{Exercises}
\begin{exercise}
Compute Poincaré duality for $S^n$.
\end{exercise}
\begin{exercise}
Compute the Betti-number symmetry for a genus-$g$ surface.
\end{exercise}
\begin{exercise}
Write the cup-product pairing matrix for $T^2$ in a standard basis.
\end{exercise}
\begin{exercise}
Identify the Poincaré dual of an embedded oriented circle in $T^2$.
\end{exercise}
\begin{exercise}
Intersect two standard generating circles on the torus.
\end{exercise}
\begin{exercise}
Relate their signed intersection to cup product.
\end{exercise}
\begin{exercise}
Compute the duality groups of $\mathbb{CP}^2$.
\end{exercise}
\begin{exercise}
Evaluate the square of the degree-two generator on $[\mathbb{CP}^2]$.
\end{exercise}
\begin{exercise}
Show that closed oriented odd-dimensional manifolds have Euler characteristic zero.
\end{exercise}
\begin{exercise}
Describe the local orientation classes defining $[M]$.
\end{exercise}
\begin{exercise}
Explain the cap-product degree shift.
\end{exercise}
\begin{exercise}
Verify naturality of evaluation under an orientation-preserving homeomorphism.
\end{exercise}
\begin{exercise}
Use a tubular neighborhood to build the dual class of a submanifold.
\end{exercise}
\begin{exercise}
Relate that construction to the Thom collapse.
\end{exercise}
\begin{exercise}
Show that a transverse zero set represents the dual of an Euler class.
\end{exercise}
\begin{exercise}
State Poincaré--Lefschetz duality for a compact manifold with boundary.
\end{exercise}
\begin{exercise}
Compute it for an interval.
\end{exercise}
\begin{exercise}
Compute it for a disk.
\end{exercise}
\begin{exercise}
Explain why ordinary integral duality fails without orientation.
\end{exercise}
\begin{exercise}
Repair the statement using mod-$2$ coefficients.
\end{exercise}
\begin{exercise}
Describe the complementary-degree pairing in dimension four.
\end{exercise}
\begin{exercise}
Identify the middle-dimensional bilinear form that will appear in VIII/34.
\end{exercise}
\begin{exercise}
Explain how transverse intersections produce its coefficients.
\end{exercise}
\begin{exercise}
Summarize fundamental class $\to$ cap product $\to$ duality $\to$ intersection form.
\end{exercise}

\section{Hints}
\begin{hint}
Only degrees $0$ and $n$ survive.
\end{hint}
\begin{hint}
Use $b_0=b_2=1$ and $b_1=2g$.
\end{hint}
\begin{hint}
Choose generators $a,b$ with $a\smile b=[T^2]^*$.
\end{hint}
\begin{hint}
Use a transverse complementary circle.
\end{hint}
\begin{hint}
There is one signed intersection.
\end{hint}
\begin{hint}
Evaluate the cup product on $[T^2]$.
\end{hint}
\begin{hint}
Use one generator in degrees $0,2,4$.
\end{hint}
\begin{hint}
The square evaluates to $1$ in the complex orientation.
\end{hint}
\begin{hint}
Pair complementary Betti numbers.
\end{hint}
\begin{hint}
Use $H_n(M,M\setminus\{x\})$.
\end{hint}
\begin{hint}
Cap lowers degree by $k$.
\end{hint}
\begin{hint}
Track the fundamental class.
\end{hint}
\begin{hint}
Use the normal Thom class.
\end{hint}
\begin{hint}
Collapse the complement of the tubular neighborhood.
\end{hint}
\begin{hint}
Use the section's zero locus.
\end{hint}
\begin{hint}
Replace absolute homology by relative homology.
\end{hint}
\begin{hint}
Check endpoints as boundary.
\end{hint}
\begin{hint}
Use $(D^n,S^{n-1})$.
\end{hint}
\begin{hint}
The local signs cannot be chosen coherently.
\end{hint}
\begin{hint}
Signs disappear modulo two.
\end{hint}
\begin{hint}
Use degree two against degree two.
\end{hint}
\begin{hint}
Evaluate cup products on $[M]$.
\end{hint}
\begin{hint}
Represent classes by transverse surfaces.
\end{hint}
\begin{hint}
Track the middle degree.
\end{hint}
